\section{工程落地蓝图：从课堂范式到可复现实验}

\subsection{先定义“可优化对象”，再优化模型}
这节课最值得迁移到日常研发的一点是：\textbf{先把任务定义成可执行、可验证、可迭代的对象，再谈模型优化}。很多团队失败并不是模型不够强，而是任务定义不完整，例如只有自然语言目标、没有明确 verifier，或只有最终分数、没有轨迹级诊断。

\begin{importantbox}{安全 Code Agent 的最小实验单元}
一个可复现实验至少应包含六个组件：
\begin{enumerate}
    \item 任务规范：输入、输出、约束、超时；
    \item 执行环境：可重建容器与依赖版本；
    \item 工具接口：代码浏览、运行、调试、输入生成；
    \item 验证器：测试集或 sanitizer/oracle；
    \item 轨迹日志：每轮 thought/tool/observation；
    \item 评测脚本：批量运行与统计汇总。
\end{enumerate}
\end{importantbox}

若缺失第 5 项和第 6 项，系统往往只能 ``看到结果、看不到过程''，导致问题复现困难。Sutton 讲 ``always look at the data'' 的真正落点就是这里：团队要能把失败轨迹当作一等数据资产持续分析。

\begin{knowledgebox}{从课堂到工程的桥接策略}
可以把课程中的三段结构映射成三层研发流程：
\begin{itemize}
    \item 基座层：代码任务评测（MBPP/HumanEval/SWE-Bench）用于建立基础能力；
    \item 系统层：harness、工具、控制流用于提升复杂任务稳定性；
    \item 应用层：安全任务中的动态验证与证据链用于检验真实价值。
\end{itemize}
\end{knowledgebox}

\subsection{迭代节奏：按失败类型拆解系统瓶颈}
课堂里多次提到 ``组合优化''，工程上可以进一步落地为失败分类。只要把失败原因稳定归类，调参优先级就会清晰得多，而不是盲目替换模型或堆叠 prompt。

\begin{table}[H]
\centering
\begin{tabular}{p{3.2cm}p{4.6cm}p{4.6cm}}
\toprule
\textbf{失败类型} & \textbf{典型症状} & \textbf{优先修复方向} \\
\midrule
检索失败 & 反复浏览无关文件，定位路径漂移 & 提升符号级检索、加入跨引用导航 \\
执行失败 & 命令调用格式错、环境依赖缺失 & 工具参数约束、容器基线固化 \\
验证失败 & 测试通过但语义错误，或误判崩溃原因 & 强化 verifier、增加反例回放 \\
策略失败 & 长轨迹自我强化错误假设 & 控制流重置、阶段化循环、状态机约束 \\
成本失败 & 成功率提升但 token/时延不可接受 & 预算管理、早停策略、分层模型路由 \\
\bottomrule
\end{tabular}
\caption{将 Agent 失败模式映射到可操作的系统改进方向}
\end{table}

\begin{warningbox}{只做“成功案例复盘”会误导团队}
如果评审会议只展示 best trajectory，系统会被优化成 ``偶尔很惊艳''。要提升线上可靠性，必须固定审查长尾失败样本，尤其是 ``看起来快成功但最终崩坏'' 的轨迹。
\end{warningbox}

\subsection{控制流治理：探索效率与可审计性的平衡}
从课程给出的连续谱来看，控制流不是选边站，而是分场景配置。探索阶段可偏动态，收敛阶段可偏约束。对安全任务而言，还需额外满足审计要求：每个关键动作可追踪、每次结论可回放。

\begin{importantbox}{建议的双模式控制流}
\begin{itemize}
    \item 探索模式：高采样、多路径、宽工具权限，用于发现潜在漏洞路径；
    \item 验证模式：低采样、强约束、固定观察点，用于复现和证据固化。
\end{itemize}
两种模式共享同一任务目标，但资源策略和可解释性要求不同。
\end{importantbox}

\begin{knowledgebox}{可执行的预算管理规则}
可以把测试时预算拆成三层：
\begin{itemize}
    \item 单轨预算：限制每条轨迹最大步数与最大运行次数；
    \item 任务预算：限制每个样例的轨迹总数与总时长；
    \item 批次预算：限制整轮评测的总 GPU/CPU 小时。
\end{itemize}
当某一层触发阈值时，系统必须输出 ``为什么停止'' 的结构化原因，避免黑箱超时。
\end{knowledgebox}

\subsection{安全落地的额外约束：责任披露与双重用途风险}
Big Sleep 案例说明 ``发现漏洞'' 已经具备现实可行性，因此落地时不能只谈技术指标。安全 Agent 需要内建责任边界，包括漏洞处理流程、对外披露节奏、以及防止能力被直接转用于攻击自动化。

\begin{warningbox}{双重用途风险必须前置治理}
越是自动化的漏洞发现系统，越要把访问权限、执行沙箱、结果脱敏、披露审批做成默认流程，而不是项目后期补丁。否则技术成功会直接转化为合规风险。
\end{warningbox}

\begin{importantbox}{实践中的最小安全治理清单}
\begin{enumerate}
    \item 所有目标仓库与运行环境白名单化；
    \item 默认只在隔离执行环境运行 exploit-like 输入；
    \item 发现疑似漏洞后先内部验证再进入披露流程；
    \item 输出报告区分 ``证据'' 与 ``推断''，避免夸大结论；
    \item 对外发布时去除可直接复现攻击链的敏感细节。
\end{enumerate}
\end{importantbox}

\subsection{研究机会：下一步值得投入的问题}
Sutton 结尾强调这是 ``wide open area''，如果转成研究议程，可以落在四类问题：更强动态 verifier、更稳定跨仓库泛化、更低成本长轨迹搜索、更可信的自动证据生成。它们共同决定安全 Agent 能否从单点成功走向规模化稳定产出。

\begin{table}[H]
\centering
\begin{tabular}{p{3.0cm}p{4.4cm}p{4.8cm}}
\toprule
\textbf{问题域} & \textbf{当前瓶颈} & \textbf{潜在突破方向} \\
\midrule
动态验证器 & 触发信号不稳定、误报判定成本高 & 组合 sanitizer + 语义断言 + 差分执行 \\
跨项目泛化 & 工具和提示高度依赖单仓库习惯 & 可迁移 ACI 协议与任务归一化接口 \\
搜索效率 & 长轨迹成本高且失败回报稀疏 & 分层规划、经验回放、失败记忆复用 \\
证据自动化 & 报告可读性与可审计性不足 & 结构化漏洞证据模板与可回放 artifact \\
\bottomrule
\end{tabular}
\caption{安全 Agent 的中短期研究机会图谱}
\end{table}

\subsection{本章小结}
把课堂方法迁移到工程实践的关键是两件事：第一，任何优化都要围绕可执行闭环；第二，任何能力都要附带审计与治理机制。只有同时满足 ``能做'' 与 ``可控''，安全 Agent 才能持续进入真实生产流程。

